%==============================================================================
\chapter{Plan de tests et validation}
%==============================================================================
La stratégie de test couvre les fonctions métier, les contraintes de gestion et
les exigences techniques. Chaque cas de test précise la fonction visée, les
préconditions, les étapes et le résultat attendu.

\section{Cas de tests}
\begin{longtable}{>{\bfseries}p{1.4cm} p{3.6cm} p{4.2cm} p{4.4cm}}
\toprule
\textbf{ID} & \textbf{Fonction} & \textbf{Étapes} & \textbf{Résultat attendu} \\
\midrule
\endhead
TC01 & CRUD ruche & Créer, lire, modifier, supprimer une ruche & Opérations conformes et persistées \\
TC02 & Contrainte cadres & Ajouter un 11e cadre à un corps & Refus, maximum de 10 cadres \\
TC03 & Contrainte hausses & Ajouter une 6e hausse & Refus, maximum de 5 hausses \\
TC04 & Contrainte socle & Retirer le dernier cadre du corps & Refus, au moins 1 cadre \\
TC05 & Planning de visite & Proposer puis approuver une visite & Statut approuvé enregistré \\
TC06 & Rapport de visite & Remplir tous les champs et enregistrer & Rapport persisté, case exécutée \\
TC07 & Alerte production & Dépasser le seuil de production & Alerte de collecte affichée \\
TC08 & Alerte baisse & Simuler une chute d'effectif & Alerte d'inspection affichée \\
TC09 & Calcul ROI & Saisir coûts et production & ROI calculé correctement \\
TC10 & API miel & Appeler getZummHoneyActualQuantity & Quantité correcte retournée \\
TC11 & Internationalisation & Basculer FR, EN puis AR & Textes traduits via le fichier XML \\
TC12 & Configuration & Modifier un seuil dans ConfigZumm.ini & Nouveau seuil pris en compte \\
TC13 & Authentification & Se connecter par OAuth Google & Accès autorisé \\
TC14 & Sécurité & Vérifier SSL et certificat X.509 & Échanges chiffrés et validés \\
TC15 & Export & Exporter les enregistrements & Fichiers CSV et TXT générés \\
TC16 & Mode hors-ligne & Saisir sans réseau puis reconnecter & Synchronisation réussie \\
\bottomrule
\end{longtable}

\section{Matrice de traçabilité}
La matrice relie chaque exigence à son cas d'utilisation, aux classes et tables
qui la portent, à l'écran (maquette) concerné et aux cas de test qui la valident.
Elle assure la traçabilité de bout en bout, du besoin jusqu'à la vérification.
{\small
\begin{longtable}{>{\bfseries}p{3.4cm} C{1.3cm} p{3.3cm} p{2.5cm} p{2.6cm}}
\toprule
\textbf{Exigence} & \textbf{UC} & \textbf{Classes / Tables} & \textbf{Écran} & \textbf{Tests} \\
\midrule
\endhead
Gestion des entités et contraintes & --- & Ruche, Compartiment, Cadre & Fiche ruche (fig.~\ref{fig:wf-ruche}) & TC01--TC04 \\
Planification et approbation & UC1 & Planning, Agent & Calendrier (fig.~\ref{fig:wf-calendrier}) & TC05 \\
Réalisation et rapport de visite & UC2 & Visite, Photo & Rapport (fig.~\ref{fig:wf-rapport}) & TC06 \\
Tableau de bord production & UC3 & Recolte, Mesure, Seuil & TB prod. (fig.~\ref{fig:wf-production}) & TC07, TC09 \\
Alerte sanitaire / baisse & UC4 & Mesure, Seuil & Calendrier & TC08 \\
Service web API (miel) & UC5 & Ruche, Recolte & --- & TC10 \\
Internationalisation & --- & (i18n XML) & Tous & TC11 \\
Configuration par seuils & --- & Seuil & TB prod. & TC12 \\
Authentification OAuth & --- & Agent & Connexion (fig.~\ref{fig:seq-oauth}) & TC13 \\
Sécurité des échanges & --- & (SSL / X.509) & --- & TC14 \\
Portabilité et export & --- & (export CSV/TXT) & Tous & TC15 \\
Mode hors-ligne & UC2 & Visite (file locale) & Rapport & TC16 \\
\bottomrule
\end{longtable}}

